
In the source language, Standard ML offers the programmer no direct control over
the memory layout for data structures in the generated code: consequently, it is
the compiler's responsibility to choose a memory layout that efficiently
utilizes the hardware of the target platform. However, the most ``literal''
translation of the source language to machine code (i.e. the one which preserves
exactly the data structures present in the user's original input) has
disasterous performance characteristics on modern CPUs: every function call and
every access into a (non-primitive) container incurs a cache miss at an
unpredictable address in the heap.

%% \cite rather than \citet because it's a citation for MLton

Each Standard ML compiler chooses different strategies to mitigate these costs:
MLton \cite{mlton-website} deploys a class of optimizations that it calls
``flattening,'' first introduced in \citet{flat-paper}. These strategies are also
used in MLton's parallel-friendly fork, MaPLe, introduced in \citet{mpl-orig}.

In this document, we will explore two novel ``flattening'' strategies: the
first, \texttt{PreFlatten} (building on top of the existing \texttt{Flatten}
pass from \citet{flat-paper}), specializes functions in order to enable
more-aggressive elimination of heap allocations at callsites. The second,
\texttt{ShallowFlatten} (building on top of the existing \texttt{DeepFlatten}
pass) transforms layouts of containers in order to eliminate pointers.

As we will see below, in practice, the \texttt{PreFlatten} optimizations
interact poorly with the existing \texttt{Chunkify} pass, and so fails to
materialize consistent wins in practice: most benchmark configurations show a
loss (on the order of a 1\% geomean regression across all impacted benchmarks,
generally driven by a handful of large outlier regressions). After controlling
for the effect of the \texttt{Chunkify} heuristic (by forcing all generated code
into a single chunk), the \texttt{ShallowFlatten} optimizations are
neutral-to-positive, generally demonstrating a geomean win on the order of 1-2\%
across all impacted benchmarks. This suggests that the \texttt{Chunkify}
heuristic may be over-tuned to the existing benchmark suite.

On the other hand, both \texttt{ShallowFlatten} transformations are
overwhelmingly positive in the single-threaded case, improving performance of
individual benchmarks by as much as 2x and materialzing geomean wins on the
order of as 10\% on impacted benchmarks. Multithreaded benchmarks show mixed
results that appear to indicate a bad interaction with MaPLe's garbage
collection heuristics: controlling for this effect by forcing a synchronous GC
pause at the beginning of each benchmark iteration shows individual benchmark
wins as large as 10x and a 50\% geomean win across impacted benchmarks.

\hl{TODO: should check benchmark data more closely}
